\daytitle{16}{电势能、电势与电势差}

\begin{learningbox}{本节导读}
上一节学习了电场的“力的属性”（电场强度 $E$），本节学习电场的“能的属性”——电势能、电势、电势差。这些概念和力学中的重力势能、高度差非常相似，可以用类比法理解。
\end{learningbox}

\textbf{本节读题路线：} 先判断题目从“力和位移”入手，还是从“能量和电势差”入手；静电力做功只看始末位置，匀强电场中再把位移投影到电场方向。比较电势先看电场方向和零电势约定；比较电势能再说明电荷正负和零势能位置。

\section*{一、学习目标}

\begin{enumerate}
  \item 能用 $W_{AB} = -\Delta E_p$ 说明静电力做功与电势能变化的关系。
  \item 能说出电势 $\varphi = \dfrac{E_p}{q}$ 和电势差 $U_{AB} = \varphi_A - \varphi_B$ 的定义含义。
  \item 能从匀强电场做功关系得到
  $U_{AB}=\vect E\cdot\overrightarrow{AB}=E\,AB\cos\theta$，并能正确使用
  $U_{AB}=Ed$ 的有向投影形式。
  \item 能根据等势面的性质判断电场方向和静电力做功。
  \item 能用电势差计算电场力做的功：$W_{AB} = qU_{AB}$。
\end{enumerate}
\section*{二、初高中衔接}

重力场中有“重力势能”和“高度差”——类比到电场中：
\begin{center}
\begin{tabular}{ll}
\toprule
重力场 & 静电场 \\
\midrule
重力 $mg$ & 电场力 $qE$ \\
重力势能 $mgh$ & 电势能 $E_p = q\varphi$ \\
高度位置 $h$ & 电势 $\varphi$ \\
高度差 $\Delta h$ & 电势差 $U$ \\
重力做功 $mg(h_A-h_B)$ & 电场力做功 $qU_{AB}$ \\
\bottomrule
\end{tabular}
\end{center}

这个类比只能帮助理解，不能直接套用：高度是几何量，电势是单位正电荷的电势能，且电势正负还受电荷源和零电势点影响。

\section*{三、核心概念}

\textbf{科学发展视角：} 现代电子设备通常不直接说“电场力做多少功”，而是说电压、电势差和能量。电池、芯片、传感器和粒子加速器都要用电势差来描述能量转化。后面遇到电子伏特（$\mathrm{eV}$）时，也是在描述“电荷经过电势差获得或失去的能量”。

\begin{conceptbox}{静电力做功与电势能}

\textbf{先看匀强静电场中的做功：} 设匀强静电场方向为 $x$ 轴正方向。电荷从 $A$ 到 $B$ 可以走任意曲线路径，把路径分成很多小段 $\Delta \vect{s}$。在匀强电场中，静电力 $\vect{F}=q\vect{E}$ 保持不变：
\[
\Delta W = \vect{F}\cdot \Delta \vect{s}=qE\Delta x
\]
其中 $\Delta x$ 是这一小段沿电场方向的投影。整条路径上的总功为：
\[
W_{AB}=\sum \Delta W=qE\sum \Delta x=qE(x_B-x_A)
\]
所以在匀强静电场中，静电力做功只与初末位置沿电场方向的距离有关，与路径形状无关。

\begin{center}
% TikZ 图：匀强电场中电场力做功与路径无关
% Day 16：电势能、电势与电势差
\begin{tikzpicture}[scale=0.82,>=Latex]
  \foreach \y in {-1.2,-0.6,0,0.6,1.2} {
    \draw[->,blue] (-0.35,\y) -- (4.35,\y);
  }
  \node[blue,above] at (0.15,1.2) {$\vect{E}$};

  \coordinate (A) at (0.35,-0.9);
  \coordinate (B) at (3.75,0.95);
  \filldraw[red] (A) circle (2pt) node[below left] {$A$};
  \filldraw[red] (B) circle (2pt) node[above right] {$B$};

  \draw[thick] (A) .. controls (0.95,0.68) and (2.15,-1.05) .. (B);
  \coordinate (S) at (1.5729,-0.1348);
  \coordinate (T) at (2.0348,-0.1102);
  \draw[orange!85!black,line width=2.6pt,line cap=round]
    (S) .. controls (1.7184,-0.1317) and (1.8726,-0.1287) .. (T);
  \filldraw[orange!85!black] (S) circle (1.8pt);
  \filldraw[orange!85!black] (T) circle (1.8pt);
  \node[orange!85!black,below] at (1.82,-0.52) {\small 截取这一小段};

  \draw[dashed,gray] (A) -- (0.35,1.45);
  \draw[dashed,gray] (B) -- (3.75,1.45);
  \draw[<->] (0.35,1.35) -- (3.75,1.35);
  \node[above] at (2.05,1.35) {$d=x_B-x_A$};

  \draw[gray,dashed] (S) -- (4.75,-0.55);
  \draw[gray,dashed] (T) -- (4.75,0.45);

  \begin{scope}[shift={(4.75,0)}]
    \draw[rounded corners=2pt,gray] (0,-1.25) rectangle (3.25,1.25);
    \node[font=\small] at (1.62,1.02) {放大这一小段};
    \foreach \y in {-0.65,0,0.65} {
      \draw[->,blue!55] (0.28,\y) -- (3.0,\y);
    }
    \node[blue,right] at (3.0,0.65) {$\vect{E}$};

    \coordinate (p) at (0.55,-0.55);
    \coordinate (q) at (2.45,0.45);
    \coordinate (h) at (2.45,-0.55);
    \filldraw[orange!85!black] (p) circle (1.4pt);
    \filldraw[orange!85!black] (q) circle (1.4pt);
    \draw[->,very thick,orange!85!black] (p) -- (q)
      node[midway,above left] {$\Delta \vect{s}$};
    \draw[dashed,orange!85!black] (q) -- (h);
    \draw[->,very thick,orange!85!black] (p) -- (h)
      node[midway,below] {$\Delta x$};
  \end{scope}
\end{tikzpicture}


\figcaption{匀强电场中的微元法：只有沿电场方向的位移投影参与做功}\label{fig:uniform-field-work-path}
\end{center}

\textbf{由路径无关引入电势能：} 既然静电力做功只由初末位置决定，就可以像重力势能一样定义电势能。规定：
\[
\boxed{W_{AB}=E_{pA}-E_{pB}=-\Delta E_p}
\]

这个关系也适用于一般静电场。这里先把“静电力做功与路径无关”作为静电场的基本性质：对点电荷形成的电场，可由后文的积分直接证明；多个静止点电荷形成的电场再由叠加原理推广。把路径分成微小段时，
\[
\mathrm{d}W=q\vect{E}\cdot\mathrm{d}\vect{s}
\]
只是计算每一小段做功的表达式，不能单独用“局部近似匀强”来证明总功已经与路径无关。基于这一静电场性质，才可在一般情形中定义电势能。

\textbf{电势能：} 电荷在电场中某点具有的势能。
\[
W_{AB} = -\Delta E_p = E_{pA} - E_{pB}
\]

\begin{itemize}
  \item 电场力做正功 $\to$ 电势能减少。
  \item 电场力做负功 $\to$ 电势能增加。
  \item 电势能是相对的，需要选取零势能参考点。
\end{itemize}

\end{conceptbox}

\begin{conceptbox}{电势与电势差}

\textbf{电势：} 在不改变原电场分布的前提下，在某点放入电荷量
$q\ne0$、且足够小的试探电荷。该试探电荷在此处的电势能与其电荷量之比为：
\[
\varphi = \frac{E_p}{q}
\]

\begin{itemize}
  \item 标量，单位：伏特（$\mathrm{V}$），$1\ \mathrm{V} = 1\ \mathrm{J/C}$。
  \item 对同一电场中的同一点，$E_p/q$ 与所取试探电荷的电荷量和正负无关，因此电势描述的是电场本身的性质。
  \item 电势的高低与电场线的方向有关：沿电场线方向电势降低。
  \item 电势是相对的，需要选零电势参考点（通常选大地或无穷远）。
\end{itemize}

\textbf{电势差（电压）：}
\[
U_{AB} = \varphi_A - \varphi_B
\]

\textbf{静电力做功与电势差：}
\[
\boxed{W_{AB} = qU_{AB}}
\]

\textbf{符号检查：} 先按下标写 $U_{AB}=\varphi_A-\varphi_B$，再乘电荷量 $q$ 判断做功正负；最后用 $W_{AB}=-\Delta E_p$ 判断电势能增减。负电荷不改变电势差的定义，只会改变 $qU_{AB}$ 的符号。

\textbf{匀强电场中电势差与场强的关系：} 设 $\theta$ 为
$\overrightarrow{AB}$ 与 $\vect E$ 的夹角，则
\[
\boxed{U_{AB}=\vect E\cdot\overrightarrow{AB}
=E\,AB\cos\theta}
\]
定义 $d=AB\cos\theta$ 为 $\overrightarrow{AB}$ 沿电场方向的
\textbf{有向投影}，便可写成
\[
\boxed{U_{AB}=Ed}
\]

\begin{itemize}
  \item $d>0$ 表示从 $A$ 到 $B$ 的投影与电场同向，$d<0$ 表示反向。
  \item 若题目中的 $d$ 只表示两等势面间的正距离，则应写
  $|U_{AB}|=Ed$，从而 $E=\dfrac{|U_{AB}|}{d}$。
  \item 场强的单位：$\mathrm{V/m}$ 和 $\mathrm{N/C}$ 等价。
\end{itemize}

\begin{center}
% TikZ 图：匀强电场中电势差与距离
% Day 16：电势能、电势与电势差
\begin{tikzpicture}[scale=0.9,>=Latex]
  % 平行板
  \draw[thick] (0,1.5) -- (0,-1.5);
  \node[left] at (0,0) {$+$};
  \draw[thick] (4,1.5) -- (4,-1.5);
  \node[right] at (4,0) {$-$};
  % 电场线
  \foreach \y in {-1.2,-0.4,0.4,1.2} {
    \draw[->,blue] (0.2,\y) -- (3.8,\y);
  }
  % A、B 两点
  \filldraw[red] (1,0.5) circle (2pt);
  \node[above] at (1,0.5) {$A$};
  \filldraw[red] (3,-0.5) circle (2pt);
  \node[below] at (3,-0.5) {$B$};
  % 距离 d
  \draw[<->] (1,0.8) -- (3,0.8);
  \node[above] at (2,0.8) {$d$};
  % 公式
  \node at (2,-1.8) {$U_{AB} = Ed$，$W_{AB} = qU_{AB}$};
\end{tikzpicture}


\figcaption{匀强电场中 $U_{AB}=Ed$：图示情形中 $d>0$}\label{fig:potential-diff}
\end{center}

\end{conceptbox}

\begin{conceptbox}{等势面}

\textbf{定义：} 电场中电势相等的点构成的曲面。

\textbf{性质：}
\begin{itemize}
  \item 等势面与电场线处处垂直。
  \item 沿等势面移动电荷，静电力不做功。
  \item 若相邻等势面取相同的电势差，等势面越密的地方场强越大（与电场线疏密的比较条件相同）。
  \item 不同等势面不相交。
\end{itemize}

\textbf{静电平衡导体：} 导体内部场强为零，导体是等势体，导体表面也是等势面；紧邻导体表面的电场方向与表面垂直。这里的“等势”不表示电势一定为零，其数值仍由零电势参考点决定。

\begin{center}
% TikZ 图：等势面与电场线
% Day 16：电势能、电势与电势差
\begin{tikzpicture}[scale=0.8,>=Latex]
  % 正电荷
  \filldraw[red] (0,0) circle (3pt);
  \node[above right] at (0,0) {$+Q$};
  % 电场线（含结束箭头）
  \foreach \angle in {0,45,90,135,180,225,270,315} {
    \draw[->,thick] (0,0) -- (\angle:1.2);
    \draw[->,thick] (\angle:1.2) -- (\angle:2.4);
  }
  % 等势面（虚线圆）
  \draw[dashed,thick] (0,0) circle (1.0);
  \draw[dashed,thick] (0,0) circle (1.8);
  \draw[dashed,thick] (0,0) circle (2.6);
  \node[below right] at (1.8,0) {等势面};
\end{tikzpicture}


\figcaption{点电荷与静电平衡导体的等势特征：左图是点电荷三维等势球面的平面截线}\label{fig:equipotential}
\end{center}

\end{conceptbox}

\begin{samepage}
\section*{四、公式与适用条件}

\textbf{公式使用前检查：}
\nopagebreak[4]
\begin{enumerate}
  \item 先规定零电势点或只讨论电势差，避免把“电势”当成绝对量。
  \item 电场力做功先看电荷符号和位移方向，再判断电势能增减。
  \item 匀强电场中用 $U_{AB}=Ed$ 时，$d$ 是
  $\overrightarrow{AB}$ 沿电场方向的有向投影。
  \item 点电荷电势是标量叠加，但电场强度仍是矢量叠加。
\end{enumerate}
\end{samepage}

\begin{center}
\begin{tabular}{lll}
\toprule
公式 & 物理意义 & 适用条件 \\
\midrule
$W_{AB} = qU_{AB}$ & 静电力做功 & 静电场 \\
$\varphi = \dfrac{E_p}{q}$ & 电势定义 & 静电场 \\
$U_{AB} = \varphi_A - \varphi_B$ & 电势差 & 静电场 \\
$U_{AB}=E\,AB\cos\theta=Ed$ & 场强与电势差 & 匀强静电场，$d$ 为有向投影 \\
沿电场线方向电势降低 & 电势变化规律 & 静电场 \\
\bottomrule
\end{tabular}
\end{center}

\begin{learningbox}{电势题的符号判断清单}
电势题的正负号应按定义逐步确定：
\begin{enumerate}
  \item \textbf{先写下标：} $U_{AB}=\varphi_A-\varphi_B$，表示从 $A$ 到 $B$ 的电势差。下标一换，符号就变：$U_{BA}=-U_{AB}$。
  \item \textbf{再乘电荷：} $W_{AB}=qU_{AB}$。正电荷和负电荷经过同一电势差，电场力做功符号可能相反。
  \item \textbf{最后看能量：} $\Delta E_p=E_{pB}-E_{pA}=-W_{AB}$。电场力做正功，电势能减少；电场力做负功，电势能增加。
\end{enumerate}
\textbf{读题提醒：} 比较电势只看电场本身；比较电势能必须同时看电荷正负。负电荷在高电势处的电势能反而可能更低。
\end{learningbox}

\section*{五、典型模型}

\subsection*{模型 1：匀强电场中的电势差}

平行板间电势沿电场线方向均匀降低。若
$d=\overrightarrow{AB}$ 沿电场方向的有向投影，则
$U_{AB}=Ed$；若只给出正的几何间距，则使用 $|U_{AB}|=Ed$。

\begin{center}
% TikZ 图：匀强电场中的电势差与投影距离
% Day 16：电势能、电势与电势差
\begin{tikzpicture}[scale=0.88,>=Latex]
  % 平行板与电场线
  \draw[thick] (-0.2,1.5) -- (-0.2,-1.5);
  \node[left] at (-0.2,0) {$+$};
  \draw[thick] (4.4,1.5) -- (4.4,-1.5);
  \node[right] at (4.4,0) {$-$};
  \foreach \y in {-1.0,-0.35,0.35,1.0} {
    \draw[->,blue] (0.05,\y) -- (4.15,\y);
  }
  \node[blue] at (4.45,1.05) {$\vect E$};

  % 等势面
  \foreach \x/\lab in {0.8/高电势,2.2/,3.6/低电势} {
    \draw[gray,dashed] (\x,1.35) -- (\x,-1.35);
  }
  \node[gray,font=\small] at (0.8,1.65) {高电势};
  \node[gray,font=\small] at (3.6,1.65) {低电势};

  % A、B 与投影
  \coordinate (A) at (0.8,0.75);
  \coordinate (B) at (3.6,-0.55);
  \filldraw[red] (A) circle (2pt) node[above left] {$A$};
  \filldraw[red] (B) circle (2pt) node[below right] {$B$};
  \draw[red,thick] (A) -- (B);
  \draw[dashed,gray] (A) -- (0.8,-1.22);
  \draw[dashed,gray] (B) -- (3.6,-1.22);
  \draw[<->] (0.8,-1.22) -- (3.6,-1.22);
  \node[below] at (2.2,-1.22) {$d$};
  \node[font=\small] at (2.2,-1.65) {$d$ 是 $AB$ 沿电场方向的投影};
\end{tikzpicture}


\figcaption{匀强电场中的电势差：只看沿电场方向的投影距离}\label{fig:uniform-field-potential-gradient}
\end{center}

若 $A$、$B$ 连线不沿电场方向，不能直接把 $AB$ 当作
$d$ 代入 $U_{AB}=Ed$，必须先求有向投影。

\subsection*{模型 2：点电荷的电势}

点电荷 $Q$ 周围距 $r$ 处的电势（取无穷远为零势点）：
\[
\varphi = k\frac{Q}{r}
\]
正电荷周围电势为正，负电荷周围电势为负。

\textbf{推导：} 以正点电荷为例，把正试探电荷 $q$ 从距 $Q$ 为 $r$ 的点移到无穷远。积分号
$\int$ 表示把连续路径上无数个微小功 $\mathrm dW$ 相加。用
$\rho$ 表示积分过程中的径向位置，则电场力做功为
\[
W_{r\to\infty}=\int_r^\infty k\frac{Qq}{\rho^2}\,\mathrm{d}\rho
=kQq\left[-\frac{1}{\rho}\right]_r^\infty
=k\frac{Qq}{r}
\]
取无穷远处电势能为零，则 $W_{r\to\infty}=E_p(r)-E_p(\infty)=E_p(r)$，所以
\[
E_p(r)=k\frac{Qq}{r},\qquad
\varphi=\frac{E_p}{q}=k\frac{Q}{r}
\]
若 $Q$ 为负，公式仍成立，电势为负。

\begin{center}
% TikZ 图：点电荷电势公式推导
% Day 16：电势能、电势与电势差
\begin{tikzpicture}[scale=0.9,>=Latex]
  \filldraw[red] (0,0) circle (4pt);
  \node[below] at (0,0) {$+Q$};
  \draw[->,thick] (0,0) -- (5.2,0);
  \node[right] at (5.2,0) {无穷远};

  \filldraw[black] (2.0,0) circle (2pt);
  \node[above left=2pt] at (2.0,0) {$P(r)$};
  \draw[<->] (0,-0.55) -- (2.0,-0.55);
  \node[below] at (1.0,-0.55) {$r$};

  \draw[->,blue,thick] (2.15,0.45) -- (3.15,0.45);
  \node[blue,above] at (2.65,0.45) {$\vect{E}$};
  \node[font=\small] at (4.05,0.75) {$E=k\dfrac{Q}{r^2}$};

  \draw[->,orange!80!black,thick] (2.0,-0.95) -- (3.05,-0.95);
  \node[orange!80!black,below] at (2.52,-0.95) {$\mathrm{d}r$};
\end{tikzpicture}


\figcaption{点电荷电势公式的推导：取无穷远为零电势点}\label{fig:point-charge-potential-derivation}
\end{center}

\textbf{多个点电荷的电势叠加：} 在同一零电势参考下，各点电荷在
$P$ 点产生的电势作代数相加：
\[
\boxed{\varphi(P)=\sum_i k\frac{Q_i}{r_i}}
\]
其中 $Q_i$ 按正、负号代入，$r_i$ 是 $P$ 到第 $i$ 个点电荷的距离。电势是标量，因此不需要按方向分解；但场强仍须作矢量叠加。

\begin{center}
% TikZ 图：多个点电荷电势的标量叠加
\begin{tikzpicture}[scale=0.9,>=Latex,font=\small]
  \coordinate (Qone) at (-2.15,0);
  \coordinate (Qtwo) at (2.15,0);
  \coordinate (P) at (0.45,1.85);

  \filldraw[red] (Qone) circle (3.5pt);
  \node[below=3pt] at (Qone) {$+Q_1$};
  \filldraw[blue] (Qtwo) circle (3.5pt);
  \node[below=3pt] at (Qtwo) {$Q_2<0$};
  \filldraw (P) circle (2.3pt);
  \node[above=2pt] at (P) {$P$};

  \draw[dashed,gray] (Qone) -- (P)
    node[midway,above left,fill=white,inner sep=1pt] {$r_1$};
  \draw[dashed,gray] (Qtwo) -- (P)
    node[midway,above right,fill=white,inner sep=1pt] {$r_2$};

  \node at (0,-1.15)
    {$\displaystyle \varphi(P)=k\frac{Q_1}{r_1}+k\frac{Q_2}{r_2}$};
\end{tikzpicture}


\figcaption{多个点电荷的电势作标量叠加，正、负号随源电荷代入}\label{fig:potential-superposition}
\end{center}

\subsection*{模型 3：带电粒子在电场中的能量关系}

只有静电力做功时，由动能定理
\[
\Delta E_k=W_{AB}=qU_{AB}
\]
以及 $W_{AB}=-\Delta E_p$，可得
\[
\Delta(E_k+E_p)=0
\]
即带电粒子的动能和电势能之和守恒：
\[
\frac12 mv_1^2 + q\varphi_1 = \frac12 mv_2^2 + q\varphi_2
\]
写成初、末速度关系为
\[
\boxed{\frac12mv_B^2-\frac12mv_A^2=qU_{AB}}
\]
正电荷从高电势向低电势运动时，电势能减少、动能增加；负电荷则要结合 $qU_{AB}$ 的符号判断，不能只凭运动方向下结论。

\begin{center}
% TikZ 图：带正电粒子在匀强电场中的能量转化
\begin{tikzpicture}[scale=0.88,>=Latex,font=\small]
  \draw[thick] (-2.7,1.35) -- (-2.7,-1.25);
  \node[left] at (-2.7,0) {$+$};
  \draw[thick] (2.7,1.35) -- (2.7,-1.25);
  \node[right] at (2.7,0) {$-$};
  \foreach \y in {-0.85,-0.28,0.28,0.85} {
    \draw[->,blue] (-2.45,\y) -- (2.45,\y);
  }
  \node[blue,above] at (0,0.85) {$\vect E$};

  \filldraw[red] (-1.55,0) circle (3pt);
  \node[red,above=3pt] at (-1.55,0) {$+q,\ A$};
  \draw[->,red,thick] (-1.40,-0.42) -- (-0.72,-0.42)
    node[midway,below] {$v_A$};

  \filldraw[red] (1.10,0) circle (3pt);
  \node[red,above=3pt] at (1.10,0) {$B$};
  \draw[->,red,thick] (1.25,-0.42) -- (2.22,-0.42)
    node[midway,below] {$v_B$};

  \node[align=center] at (0,-1.72)
    {$U_{AB}>0,\quad qU_{AB}=\Delta E_k$\\
     $E_p\text{ 减少}\quad\longrightarrow\quad E_k\text{ 增加}$};
\end{tikzpicture}


\figcaption{正电荷在匀强电场中由高电势向低电势运动：电势能转化为动能}\label{fig:electric-potential-energy-conversion}
\end{center}

\textbf{电子伏特：} 一个元电荷通过 $1\ \mathrm V$ 的电势差时，能量变化的大小定义为
\[
1\ \mathrm{eV}=e\times1\ \mathrm V=1.602\times10^{-19}\ \mathrm J.
\]
$\mathrm{eV}$ 是能量单位，不是电压单位。

\newpage
\section*{六、例题精讲}

\begin{examplebox}{电场力做功}

\textbf{题意分析：} 在匀强电场中，$B$ 位于 $A$ 沿电场方向的一侧。将电荷 $q = +2 \times 10^{-9}\ \mathrm{C}$ 从 $A$ 点移到 $B$ 点，两点沿电场方向的距离为 $d = 0.05\ \mathrm{m}$，电场强度 $E = 200\ \mathrm{N/C}$。求：
\begin{enumerate}
  \item $A$、$B$ 间的电势差。
  \item 电场力做的功。
  \item 若 $q$ 改为 $-2 \times 10^{-9}\ \mathrm{C}$，结果如何？
\end{enumerate}

\textbf{物理模型：} 图示情形中 $d>0$，故 $U_{AB}=Ed$，再用 $W_{AB}=qU_{AB}$。

\textbf{解题过程：}

\begin{enumerate}
  \item $U_{AB} = Ed = 200 \times 0.05 = 10\ \mathrm{V}$

  \item $W_{AB} = qU_{AB} = 2 \times 10^{-9} \times 10 = 2 \times 10^{-8}\ \mathrm{J}$

  \item $q' = -2 \times 10^{-9}\ \mathrm{C}$：
  \[
  W_{AB}' = q'U_{AB} = -2 \times 10^{-9} \times 10 = -2 \times 10^{-8}\ \mathrm{J}
  \]
  负功表示电场力做负功，电势能增加。
\end{enumerate}

\textbf{方法总结：} 静电力做功题先确定初末位置和电势差，再代入 $W=qU$；路径形状不影响结果，电荷正负决定做功符号。

\end{examplebox}

\begin{examplebox}{电势的计算}

\textbf{题意分析：} 点电荷 $Q = +3 \times 10^{-9}\ \mathrm{C}$，取无穷远处电势为零。求距 $Q$ 分别为 $0.1\ \mathrm{m}$ 和 $0.3\ \mathrm{m}$ 处的电势。

\textbf{物理模型：} 点电荷的电势公式。

\textbf{解题过程：}

\[
\varphi_1 = k\frac{Q}{r_1} = 9 \times 10^9 \times \frac{3 \times 10^{-9}}{0.1} = 270\ \mathrm{V}
\]
\[
\varphi_2 = k\frac{Q}{r_2} = 9 \times 10^9 \times \frac{3 \times 10^{-9}}{0.3} = 90\ \mathrm{V}
\]

$U_{12} = \varphi_1 - \varphi_2 = 270 - 90 = 180\ \mathrm{V}$。

\textbf{方法总结：} 取无穷远为零电势点时，正点电荷周围电势为正，离电荷越近电势越高。两点间的电势差等于电势之差。

\end{examplebox}

\begin{examplebox}{静电力做功与路径无关}

\textbf{题意分析：} 在匀强电场中，将 $q = +1 \times 10^{-8}\ \mathrm{C}$ 的电荷从 $A$ 沿直线移到 $B$，电场力做功 $W_1 = 3 \times 10^{-6}\ \mathrm{J}$。若从 $A$ 沿另一条路径（先向右再向上）移到 $B$，电场力做功 $W_2$ 是多少？$A$、$B$ 的电势差是多少？

\textbf{物理模型：} 静电力做功与路径无关。

\textbf{解题过程：}

$W_2 = W_1 = 3 \times 10^{-6}\ \mathrm{J}$（与路径无关）。

\[
U_{AB} = \frac{W_{AB}}{q} = \frac{3 \times 10^{-6}}{1 \times 10^{-8}} = 300\ \mathrm{V}
\]

\textbf{方法总结：} 遇到路径很多的题，先判断是否是静电场；静电力做功只看初末位置，因此可以引入电势能。

\end{examplebox}

\begin{examplebox}{带电粒子经电势差加速}
\textbf{题意分析：} 质子从静止出发，在真空中从高电势的 $A$ 点运动到低电势的 $B$ 点，$U_{AB}=500\ \mathrm V$。取
$e=1.60\times10^{-19}\ \mathrm C$，
$m_p=1.67\times10^{-27}\ \mathrm{kg}$，求质子到达 $B$ 点时的动能和速度。

\textbf{物理模型：} 只有静电力做功，$\Delta E_k=qU_{AB}$。

\textbf{解题过程：}
\[
E_{kB}=eU_{AB}=500\ \mathrm{eV}
=8.00\times10^{-17}\ \mathrm J
\]
\[
v_B=\sqrt{\frac{2eU_{AB}}{m_p}}
\approx3.10\times10^5\ \mathrm{m/s}.
\]

\textbf{方法总结：} 求能量时可先用电子伏特表示，再按需要换算成焦耳；求速度时必须使用国际单位制。
\end{examplebox}

\section*{七、变式训练}

\begin{enumerate}
  \item 在匀强电场中将 $q = -1 \times 10^{-8}\ \mathrm{C}$ 从 $A$ 移到 $B$，$U_{AB} = 50\ \mathrm{V}$。求电场力做功。$q$ 的电势能如何变化？

  \item 点电荷 $Q = -4 \times 10^{-9}\ \mathrm{C}$，求距 $Q$ 为 $0.2\ \mathrm{m}$ 处的电势。

  \item 匀强电场中，$A$、$B$ 相距 $0.1\ \mathrm{m}$，连线与电场方向夹角 $60^\circ$，$U_{AB} = 20\ \mathrm{V}$。求电场强度 $E$。
  \begin{center}
  \begin{tikzpicture}[scale=0.78,>=Stealth]
  \foreach \y in {0,0.7,1.4,2.1} {
    \draw[->,blue] (0,\y) -- (4.8,\y);
  }
  \coordinate (A) at (1.3,0.28);
  \coordinate (B) at (2.3,2.01);
  \fill (A) circle (1.8pt) node[below left] {$A$};
  \fill (B) circle (1.8pt) node[above] {$B$};
  \draw[thick] (A) -- (B);
  \node[left] at (1.18,1.15) {$AB=0.1\ \mathrm{m}$};
  \draw[dashed] (A) -- ++(1.45,0);
  \draw[<->,orange,thick] (1.3,-0.18) -- (2.8,-0.18) node[midway,above] {$d$};
  \draw (1.88,0.28) arc[start angle=0,end angle=60,radius=0.58];
  \node at (2.12,0.72) {$60^\circ$};
  \node[right,blue] at (4.8,2.1) {$E$};
  \node[font=\small] at (2.45,-0.92) {沿电场方向的有向投影参与 $U_{AB}=Ed$};
\end{tikzpicture}

  \end{center}

  \item 两个点电荷 $Q_1=+4.0\times10^{-9}\ \mathrm C$、
  $Q_2=-2.0\times10^{-9}\ \mathrm C$。点 $P$ 到它们的距离分别为
  $0.20\ \mathrm m$ 和 $0.10\ \mathrm m$，取无穷远为零电势点。求
  $P$ 点电势，并判断能否由此断定 $P$ 点场强为零。

  \item 孤立导体达到静电平衡。比较导体内部两点 $M,N$ 与表面一点
  $S$ 的电势；将电荷 $q$ 沿导体表面从 $S$ 移到另一点，静电力做功是多少？

  \item 正电荷 $q$ 从静止出发，只在静电力作用下由 $A$ 运动到
  $B$，且 $U_{AB}>0$。写出末速度 $v_B$ 的表达式。
\end{enumerate}

\begin{learningbox}{变式训练参考结果}
\begin{enumerate}[itemsep=0.45em,topsep=0.3em,parsep=0pt]
  \item $W_{AB}=qU_{AB}=-5.0\times10^{-7}\ \mathrm J$；由
  $\Delta E_p=-W_{AB}$，电势能增加 $5.0\times10^{-7}\ \mathrm J$。
  \item 取无穷远处电势为零，$\varphi=kQ/r=-180\ \mathrm V$。
  \item $U_{AB}=E\,AB\cos60^\circ$，故 $E=400\ \mathrm{V/m}$。
  \item $\varphi_P=kQ_1/r_1+kQ_2/r_2=0$；电势为零只是标量代数和为零，不能据此断定矢量场强为零。
  \item $\varphi_M=\varphi_N=\varphi_S$；沿表面移动时电势差为零，故静电力做功为零。
  \item 由 $\frac12mv_B^2=qU_{AB}$，得
  $v_B=\sqrt{\dfrac{2qU_{AB}}{m}}$。
\end{enumerate}
\end{learningbox}

\section*{八、竞赛入门补充}

下面把“静电力做功与路径无关”从匀强电场推广到非匀强电场。关键仍是投影思想：真正参与做功的是沿电场方向的位移分量，与电场垂直的分量不做功。

\begin{competitionbox}
\textbf{非匀强电场中路径无关的证明}

匀强电场中可以直接用投影证明路径无关。非匀强电场里，场强大小和方向随位置变化，不能再简单写成 $qE d$，但静电场仍然具有路径无关性。

\textbf{以点电荷电场为例：} 点电荷 $Q$ 周围的电场是非匀强电场。电场强度沿径向，大小为
\[
E=k\frac{|Q|}{r^2}
\]

\begin{center}
% TikZ 图：非匀强电场中路径无关的径向积分思路
% Day 16：电势能、电势与电势差
\begin{tikzpicture}[scale=0.82,>=Latex]
  \filldraw[red] (0,0) circle (4pt);
  \node[below] at (0,0) {$Q$};
  \coordinate (A) at (1.2,0.95);
  \coordinate (B) at (2.75,-0.55);
  \filldraw[black] (A) circle (2pt) node[above] {$A$};
  \filldraw[black] (B) circle (2pt) node[right] {$B$};

  \draw[thick] (A) .. controls (1.9,1.45) and (2.9,0.65) .. (B);
  \draw[thick,dashed] (A) .. controls (0.95,0.05) and (1.75,-1.1) .. (B);
  \draw[dashed,gray] (0,0) -- (A);
  \draw[dashed,gray] (0,0) -- (B);
  \node[left] at (0.62,0.48) {$r_A$};
  \node[below] at (1.35,-0.26) {$r_B$};

  \draw[->,blue] (1.4,0) -- (2.0,0) node[midway,above] {$\vect{E}$};
\end{tikzpicture}
\\[0.3em]
\figcaption{非匀强点电荷电场：不同路径的静电力做功相同}\label{fig:nonuniform-field-path-proof}
\end{center}

沿任意路径移动电荷 $q$ 时，把微小位移 $\mathrm d\vect s$ 投影到径向，得到有向分量 $\mathrm dr$；与径向垂直的分量不参与做功。若把径向向外取为正，点电荷电场的径向分量可写成
\[
E_r=k\frac{Q}{r^2},
\]
这里的 $Q$ 按带符号电荷量代入。因此微元功为
\[
\mathrm{d}W=qE_r\,\mathrm{d}r=k\frac{Qq}{r^2}\,\mathrm{d}r
\]
从 $A$ 到 $B$ 的总功为
\[
W_{AB}=\int_{r_A}^{r_B} k\frac{Qq}{r^2}\,\mathrm{d}r
=kQq\left(\frac{1}{r_A}-\frac{1}{r_B}\right)
\]
结果只含 $r_A$、$r_B$，与具体路径无关。

\textbf{推广：} 多个静止点电荷形成的静电场可以看成许多点电荷电场的叠加。每一部分做功都只与初末位置有关，所以总功也只与初末位置有关。因此静电场可以定义电势能和电势。
\end{competitionbox}

\begin{competitionbox}
\textbf{匀强电场中电势的几何关系}

在匀强电场中，沿任一固定直线方向，电势随位移在电场方向上的有向投影线性变化；若该方向与电场垂直，电势保持不变。这一性质是处理几何型电势问题的重要依据。

\textbf{结论：} 在匀强电场中，如果 $A$、$B$、$C$ 是同一等势面上的三个点，则它们电势相等。两个等势点只能提示等势线经过这两点；结合匀强电场中等势线为直线的性质，或结合更多等势点，才能确定相应的等势线。电场方向总是垂直等势线，并从高电势指向低电势。

\textbf{例题：} 一个边长为 $a$ 的等边三角形 $ABC$ 放在匀强电场中，已知 $\varphi_A = 10\ \mathrm{V}$、$\varphi_B = 6\ \mathrm{V}$、$\varphi_C = 2\ \mathrm{V}$。求电场强度的大小和方向。

\begin{center}
% TikZ 图：匀强电场中等边三角形电势
% Day 16：电势能、电势与电势差
\begin{tikzpicture}[scale=1.05,>=Latex]
  \coordinate (A) at (0,0);
  \coordinate (C) at (4,0);
  \coordinate (B) at (2,3.464);
  \coordinate (M) at (2,0);

  \draw[thick] (A) -- (B) -- (C) -- cycle;
  \filldraw[red] (A) circle (2.2pt) node[below left] {$A,\ 10\,\mathrm{V}$};
  \filldraw[red] (B) circle (2.2pt) node[above] {$B,\ 6\,\mathrm{V}$};
  \filldraw[red] (C) circle (2.2pt) node[below right] {$C,\ 2\,\mathrm{V}$};
  \filldraw[black] (M) circle (1.6pt) node[below right] {$M,\ 6\,\mathrm{V}$};

  % 等势线 BM 与电场方向
  \draw[gray,dashed,thick] (B) -- (M);
  \node[gray,right] at (2.05,1.75) {等势线};
  \draw[->,blue,very thick] (0.65,0.48) -- (3.35,0.48) node[midway,above] {$\vect E$};
  % 长度标注
  \draw[<->] (0,-0.85) -- (4,-0.85);
  \node[below] at (2,-0.85) {$a$};
\end{tikzpicture}


\figcaption{匀强电场中的等边三角形电势：先找等势线，再找电场方向}\label{fig:uniform-field-triangle-potential}
\end{center}

\textbf{解：} 取 $AC$ 的中点为 $M$。由于匀强电场中同一直线上的电势线性变化，且
\[
\varphi_B=\frac{\varphi_A+\varphi_C}{2}
\]
说明 $B$ 点电势等于 $AC$ 中点 $M$ 的电势，$BM$ 是等势线。电场方向垂直于等势线 $BM$，即平行于 $AC$，并由高电势的 $A$ 指向低电势的 $C$。

$A$、$C$ 间电势差为
\[
U_{AC}=\varphi_A-\varphi_C=8\ \mathrm{V}
\]
沿电场方向的距离就是 $AC$ 的长度
\[
d=a
\]
所以电场强度大小为
\[
E=\frac{U_{AC}}{d}=\frac{8\ \mathrm V}{a}
\]
若 $a$ 用米作单位，则 $E=(8/a)\ \mathrm{V/m}$。

方向：平行于 $AC$，由 $A$ 指向 $C$。

\textbf{竞赛提示：} 在匀强电场中，“两点间的电势差 = 场强 × 沿电场方向的距离”——先用等势面判断电场方向，再计算大小。
\end{competitionbox}

\newpage
\section*{九、易错点}

\begin{mistakebox}
\begin{enumerate}
  \item \textbf{$U_{AB}=Ed$ 中 $d$ 的含义：} 先把
  $\overrightarrow{AB}$ 投影到电场方向。$d$ 是有向投影，不是简单地把两点间距代入；若 $d$ 只表示正距离，应写 $|U_{AB}|=Ed$。

  \item \textbf{电势的正负：} 电势是标量，有正负。取无穷远为零电势点时，正点电荷周围电势为正，负点电荷周围电势为负。电势的大小比较时，$+5\ \mathrm{V} > -10\ \mathrm{V}$。

  \item \textbf{静电力做功正负的判断：} 先规定 $U_{AB}$ 的方向，再代 $W_{AB}=qU_{AB}$；$q$ 和 $U_{AB}$ 同号则做正功。

  \item \textbf{零电势点的选取：} 比较电势值或代入 $E_p=q\varphi$ 时，要在同一零电势参考点下进行。若只求两点电势差，零点选择会相互抵消，但仍要保持符号约定一致。

  \item \textbf{混淆电势和电势能：} 先问这个量属于“场”还是“电荷在场中”。电势 $\varphi$ 属于电场，电势能 $E_p=q\varphi$ 属于电荷。
\end{enumerate}
\end{mistakebox}

\section*{十、今日小结}

\begin{enumerate}
  \item 静电场中，静电力做功 $W_{AB}=qU_{AB}$，与路径无关。
  \item 静电力做功等于电势能变化的负值：$W_{AB}=-\Delta E_p$。
  \item 电势 $\varphi = \dfrac{E_p}{q}$，电势差 $U_{AB} = \varphi_A - \varphi_B$。
  \item 沿电场线方向电势降低。
  \item 匀强电场中
  $U_{AB}=E\,AB\cos\theta=Ed$；若 $d$ 是正距离，则
  $E=\dfrac{|U_{AB}|}{d}$。
  \item 等势面与电场线垂直，沿等势面移动电荷时静电力不做功。
  \item 多个点电荷的电势作标量代数叠加；只有静电力做功时，$E_k+E_p$ 守恒。
\end{enumerate}

\textbf{一分钟复盘：} 任选两点 $A,B$，先沿电场方向判断电势高低，再写出 $U_{AB}$、$W_{AB}=qU_{AB}$ 和 $\Delta E_p=-W_{AB}$；最后检查下标顺序和电荷正负有没有弄反。
